Negative-control estimates can be aligned in sign before pooling, but orientation alone does not establish that heterogeneous controls measure commensurable residual confounding. We retain the sign-aware equal-pooling rule and assess its conditions and limitations in educational data. A bias--variance identity gives the exact finite-sample condition for MSE improvement; it does not identify unknown target bias. In known-truth binary simulations, matched controls reduce RMSE, whereas nonmatching magnitudes and incompatible factors produce poor coverage. New common-design audits use checksum-verified OULAD inputs and an explicitly named XES3G5M sequence source. All outcomes share adjustment, cross-fit folds and support, and real-control timing and coding are documented. OULAD leave-one-control-out null prediction does not improve overall, and orientations for two consequential controls have weak marginal proxy support. XES null-reference errors decrease in development and holdout, conditional on additional observation-process assumptions. The results support conditional use and diagnostic reporting rather than a generally identified causal-calibration method. Neither non-rejection of a control null nor a target interval crossing zero validates target-bias correction.


oindent\textbf{Keywords:} negative controls; calibration; unmeasured confounding; educational data mining; sensitivity analysis.
